% !TEX TS-program = pdflatex
% !TeX program = pdflatex
% !TEX encoding = UTF-8
% !TEX spellcheck = fr

\documentclass[xcolor=table]{beamer}


%\usepackage{fullpage}
%\usepackage[left=2.8cm,right=2.2cm,top=2 cm,bottom=2 cm]{geometry}
\setbeamersize{text margin left=10pt,text margin right=10pt}


\usepackage{amsmath,amssymb} 
\usepackage{fontenc}
\usepackage{textcomp}
\usepackage[utf8]{inputenc}
\usepackage{CJKutf8}
\usepackage[french]{babel}
\usepackage{arabtex,acolor}
\usepackage{txfonts}
\usepackage{tipa}
\usepackage[]{graphicx}
\usepackage{multirow}
\usepackage{hyperref}
\usepackage{colortbl}
\usepackage{tabularray}
\usepackage{listingsutf8}
\usepackage{wrapfig}
\usepackage{multicol}
\usepackage[export]{adjustbox} %for images in table, also for frame
\usepackage[many]{tcolorbox}
\usepackage{wasysym}
\usepackage[french,lined]{algorithm2e}
\usepackage{alltt} %verbatim with commands
\usepackage{longtable}
\usepackage{tabu}


\definecolor{lightblue}{HTML}{D0D2FF}
\definecolor{lightyellow}{HTML}{FFFFAA}
\definecolor{darkblue}{HTML}{0000BB}
\definecolor{olivegreen}{HTML}{006600}
\definecolor{darkgreen}{HTML}{008B45} %009B55
\definecolor{violet}{HTML}{6600CC}
%\definecolor{deeppink}{HTML}{FF1493}
\definecolor{orangey}{HTML}{FFBB00}

\usetheme{Warsaw} % Antibes Boadilla Warsaw

\beamertemplatenavigationsymbolsempty


\let\oldcite\cite
\renewcommand{\cite}[1]{{\bfseries\color{orangey}\oldcite{#1}}}

\tcbuselibrary{listings}

\newcommand{\kurl}[1]{{\scriptsize\bfseries\color{orangey}\url{#1}}}


%\renewcommand{\baselinestretch}{1.5}

\def\supit#1{\raisebox{0.8ex}{\small\it #1}\hspace{0.05em}}

\AtBeginSection{%
	\begin{frame}
		\sectionpage
	\end{frame}
}

\newcommand{\rottext}[2]{%
	\rotatebox{90}{%
	\begin{minipage}{#1}%
		\raggedleft#2%
	\end{minipage}%
	}%
}


\institute{ %
Laboratoire de la Communication dans les Systèmes Informatiques (LCSI)

\vspace{6pt}
École  nationale Supérieure d'Informatique (ESI, ex. INI), Alger, Algérie
}
\author[ \textbf{\footnotesize\insertframenumber/\inserttotalframenumber} \hspace*{1cm} ESI - ARIES Abdelkrime (Master 2022/2023)] %
{ARIES Abdelkrime}
%\titlegraphic{\includegraphics[height=1cm]{../img/esi-logo.png}%\hspace*{4.75cm}~
	
	
	\date{Master 2022/2023} %\today

\titlegraphic{%
	\includegraphics[height=1cm]{../img/esi-logo.png}%
	\hspace{2cm}
	\includegraphics[height=1cm]{../img/lcsi-logo.png}%
}


%\setbeamertemplate{headline}{}



\newcommand{\keyword}[1]{\textcolor{red}{\bfseries\itshape #1}}
\newcommand{\expword}[1]{\textcolor{olivegreen}{#1}}
\newcommand{\optword}[1]{\textcolor{violet}{\bfseries #1}}

\makeatletter
\newcommand\mysphere{%
	\parbox[t]{10pt}{\raisebox{0.2pt}{\beamer@usesphere{item projected}{bigsphere}}}}
\makeatother

%\let\oldtabular\tabular
%\let\endoldtabular\endtabular
%\renewenvironment{tabular}{\rowcolors{2}{white}{lightblue}\oldtabular\rowcolor{blue}}{\endoldtabular}


\NoAutoSpacing %french autospacing after ":"

\def\graphpath{}

\newcommand{\changegraphpath}[1]{\def\graphpath{#1}}


\newcommand{\vgraphpage}[2][.84\textheight]{%
%	\begin{center}%
		\includegraphics[height=#1]{\graphpath #2}%
%	\end{center}%
}

\newcommand{\hgraphpage}[2][\textwidth]{%
%	\begin{center}%
		\includegraphics[width=#1]{\graphpath #2}%
%	\end{center}%
}

\newcommand{\graphpage}[2][]{%
	\includegraphics[#1]{\graphpath #2}%
}

\bibliographystyle{apalike}

\newcommand{\insertbibliography}[2]{
	\appendix
	\section*{Bibliographie}
	\nocite{#2}
%	\makeatletter % to change template
%	\setbeamertemplate{headline}[default] % not mandatory, but I though it was better to set it blank
%	\def\beamer@entrycode{\vspace*{-\headheight}} % here is the part we are interested in :)
%	\makeatother
	\begin{multicols*}{2}[\frametitle{\insertsection} \usebeamertemplate*{frametitle}]%\usebeamertemplate*{frametitle}\frametitle{Références}
		\tiny
		\bibliography{#1}
	\end{multicols*}
}

\definecolor{my-grey}{RGB}{233, 233, 233}

\newcommand{\insertlicence}{
	\begin{frame}[plain]
	\frametitle{Licence : CC-BY 4.0}
%	\framesubtitle{Licence: CC-BY-NC 4.0}

	\begin{tcolorbox}[colback=cyan,
		colframe=cyan,  
		arc=0pt,outer arc=0pt,
		valign=top, 
		halign=center,
		width=\textwidth]
		
		\includegraphics[width=.5cm]{../img/licence/cc_icon_white_x2.png}
		\includegraphics[width=.5cm]{../img/licence/attribution_icon_white_x2.png}
		
		\color{white}
		\bfseries Attribution 4.0 International (CC BY 4.0) \\
		\tiny \url{https://creativecommons.org/licenses/by/4.0/deed.fr}
		
	\end{tcolorbox}\vspace{-.5cm}
	\begin{tcolorbox}[colback=my-grey,
		colframe=my-grey,  
		center, arc=0pt,outer arc=0pt,
		valign=top, 
		halign=left,
		width=\textwidth]
		
		\tiny
		
		\begin{center}
			\bfseries\Large
			Vous êtes autorisé à :
		\end{center}
		
		\begin{minipage}{0.83\textwidth}
			\begin{itemize}
				\item[] \textbf{Partager} — copier, distribuer et communiquer le matériel par tous moyens et sous tous formats
				\item[] \textbf{Adapter} — remixer, transformer et créer à partir du matériel
				pour toute utilisation, y compris commerciale.
			\end{itemize}
		\end{minipage}
		\begin{minipage}{0.15\textwidth}
			\includegraphics[width=\textwidth]{../img/licence/FreeCulturalWorks_seal_x2.jpg}
		\end{minipage}
	
		
		\begin{center}
			\bfseries\Large
			Selon les conditions suivantes :
		\end{center}
		
		\begin{itemize}
			\item[] \textbf{Attribution} — Vous devez créditer l'Œuvre, intégrer un lien vers la licence et indiquer si des modifications ont été effectuées à l'Oeuvre. Vous devez indiquer ces informations par tous les moyens raisonnables, sans toutefois suggérer que l'Offrant vous soutient ou soutient la façon dont vous avez utilisé son Oeuvre. 
			\item[] \textbf{Pas de restrictions complémentaires} — Vous n'êtes pas autorisé à appliquer des conditions légales ou des mesures techniques qui restreindraient légalement autrui à utiliser l'Oeuvre dans les conditions décrites par la licence.
		\end{itemize}
		
	\end{tcolorbox}
	
%	\begin{center}
%		\bfseries Attribution 4.0 International (CC BY 4.0)
%		\url{https://creativecommons.org/licenses/by/4.0/deed.fr}
%	\end{center}

%	\tiny
%
%	Vous êtes autorisé à : 
%	\begin{itemize}
%		\item \textbf{Partager} — copier, distribuer et communiquer le matériel par tous moyens et sous tous formats
%		\item \textbf{Adapter} — remixer, transformer et créer à partir du matériel
%	\end{itemize}
%	
%	Selon les conditions suivantes : 
%	\begin{itemize}
%		\item \textbf{Attribution} — Vous devez créditer l'Œuvre, intégrer un lien vers la licence et indiquer si des modifications ont été effectuées à l'Oeuvre. Vous devez indiquer ces informations par tous les moyens raisonnables, sans toutefois suggérer que l'Offrant vous soutient ou soutient la façon dont vous avez utilisé son Oeuvre.
%		\item \textbf{Pas d'Utilisation Commerciale} — Vous n'êtes pas autorisé à faire un usage commercial de cette Oeuvre, tout ou partie du matériel la composant. 
%		\item \textbf{Pas de restrictions complémentaires} — Vous n'êtes pas autorisé à appliquer des conditions légales ou des mesures techniques qui restreindraient légalement autrui à utiliser l'Oeuvre dans les conditions décrites par la licence.
%	\end{itemize}

	\end{frame}
}

\settowidth{\leftmargini}{\usebeamertemplate{itemize item}}
\addtolength{\leftmargini}{\labelsep}

\AtBeginDocument{
	\newcolumntype{L}[2]{>{\vbox to #2\bgroup\vfill\flushleft}p{#1}<{\egroup}} 
	
	\begin{frame}[plain]
		\maketitle
	\end{frame}

	\insertlicence
}


% needs etoolbox; to break links after -
\appto\UrlBreaks{\do\-}


\makeatletter
\newcommand{\xRightarrow}[2][]{\ext@arrow 0359\Rightarrowfill@{#1}{#2}}
\makeatother


\usefonttheme{structurebold}
%\usefonttheme{professionalfonts}

\title[TALN : 02- Traitements basiques]%
{Traitement automatique du langage naturel\\Chapitre 02 : Traitements basiques du texte} 

\changegraphpath{../img/basique/}

\begin{document}
	
\begin{frame}
\frametitle{Traitement automatique du langage naturel}
\framesubtitle{Traitements basiques du texte : Introduction}

\begin{itemize}
	\item 
\end{itemize}
%\begin{minipage}{0.78\textwidth}
%\begin{itemize}
%	\item Un domaine multidisciplinaire
%	\begin{itemize}
%		\item \optword{Linguistique} : étude du langage
%		\item \optword{Informatique} : Traitement automatique de l'information
%		\item \optword{Intelligence artificielle} : Ensemble de théories et de techniques mises en œuvre en vue de réaliser des machines capables de simuler l'intelligence humaine.
%	\end{itemize}
%\end{itemize}
%\end{minipage}
%\begin{minipage}{0.20\textwidth}
%\hgraphpage[\textwidth]{TALN.pdf}
%\end{minipage}
\end{frame}

\begin{frame}
\frametitle{Traitement automatique du langage naturel}
\framesubtitle{Traitements basiques du texte : Un peu d'humour}

%\begin{center}
%	\vgraphpage{humour.jpg}
%\end{center}

\end{frame}

\begin{frame}
\frametitle{Traitement automatique du langage naturel}
\framesubtitle{Traitements basiques du texte : Plan}

\begin{multicols}{2}
%	\small
\tableofcontents
\end{multicols}
\end{frame}

%===================================================================================
\section{Caractères}
%===================================================================================

\begin{frame}
\frametitle{Traitements basiques du texte}
\framesubtitle{Caractères}


\end{frame}

\subsection{Expressions régulières}

\begin{frame}
\frametitle{Traitements basiques du texte : Caractères}
\framesubtitle{Expressions régulières (1)}

\begin{itemize}
	\item utilisées pour chercher des chaînes de caractères dans un texte
\end{itemize}

\rowcolors{2}{lightblue}{lightyellow}
\begin{tabular}{p{.1\textwidth}p{.34\textwidth}p{.46\textwidth}}
	\rowcolor{darkblue}
	\textcolor{white}{ER} & \textcolor{white}{Sens} & \textcolor{white}{Exemple} \\
	
	. & n'importe quel caractère & \keyword{beg.n} : I \expword{begun} at the \expword{begin}ning. \\
	
	 \empty [aeuio] & caractères spécifiques & \keyword{[Ll][ae]} : \expword{Le} chat mange \expword{la} sourie. \\
	 
	\empty [a-e] & plage de caractères & \keyword{[A-Z]..} : \expword{J'a}i vu \expword{Kar}im. \\
	
	\empty [\textasciicircum aeuio] & exclure des caractères & \keyword{[\textasciicircum A-Z]a.} : J\expword{'ai} vu Karim. \\
	
	c? & un ou zéro & \keyword{colou?r} : It is \expword{colour} or \expword{color}. \\
	
	c* & zéro ou plus & \keyword{No*n} : \expword{Nn}! \expword{Non}! \expword{Nooooooon}! \\
	
	c+ & un ou plus & \keyword{No+n} : Nn! \expword{Non}! \expword{Nooooooon}! \\
	
	c\{n\} & n occurrences & \keyword{No\{3\}n} : Nn! Non! Noon! \expword{Nooon}! \\
	
	c\{n,m\} & de n à m occurrences & \keyword{No\{1,2\}n} : Nn! \expword{Non}! \expword{Noon}! Nooon! \\
	
	c\{n,\} & au moins n occurrences & \keyword{No\{2,\}n} : Nn! Non! \expword{Noon}! \expword{Nooon}! \\
	
	c\{,m\} & au plus m occurrences & \keyword{No\{,2\}n} : \expword{Nn}! \expword{Non}! \expword{Noon}! Nooon! \\
	
\end{tabular}

\end{frame}

\begin{frame}
\frametitle{Traitements basiques du texte : Caractères}
\framesubtitle{Expressions régulières (2)}

\begin{minipage}{.6\textwidth}
\begin{itemize}
	\item pour le groupement, utiliser \keyword{( )}
	\begin{itemize}
		\item Exemple, \expword{/(bla)+/} : Ceci est du \expword{blabla}.
	\end{itemize}
	\item pour la disjonction, utiliser \keyword{\textbar}
	\begin{itemize}
		\item Exemple, \expword{/continu(er\textbar ation\textbar el(le)?s?)/}
	\end{itemize}
	\item pour le début du texte, utiliser \keyword{\textasciicircum}
	\begin{itemize}
		\item Exemple, \expword{/\textasciicircum K/} :  \expword{K}ill Karim.
	\end{itemize}
	\item pour la fin du texte, utiliser \keyword{\$}
	\begin{itemize}
		\item Exemple, \expword{/\textbackslash .[\textasciicircum .]+\$/} :  fichier.tar\expword{.gz}
	\end{itemize}
	\item pour capturer un groupe, utiliser le regroupement avec \keyword{\$n} ou \keyword{\textbackslash n} où \expword{n} est la position du groupe
	\begin{itemize}
		\item Exemple, \expword{/(.*)(ement\textbar ation)\$/\textbackslash 1er/}
	\end{itemize}
\end{itemize}
\end{minipage}
\begin{minipage}{.38\textwidth}
\rowcolors{2}{lightblue}{lightyellow}
\begin{tabular}{p{.2\textwidth}p{.6\textwidth}}
	\rowcolor{darkblue}
	\textcolor{white}{ER} & \textcolor{white}{Équivalence} \\
	
	\textbackslash d & [0-9] \\
	\textbackslash D & [\textasciicircum 0-9] \\
	\textbackslash w & [a-zA-Z0-9\_] \\
	\textbackslash W & [\textasciicircum \textbackslash w] \\
	\textbackslash s & [ \textbackslash r\textbackslash t\textbackslash n\textbackslash f] \\
	\textbackslash S & [\textasciicircum \textbackslash s] \\
\end{tabular}

\end{minipage}

\end{frame}

\begin{frame}
\frametitle{Traitements basiques du texte : Caractères}
\framesubtitle{Expressions régulières : Cas d'utilisation}

\begin{itemize}
	\item Les éditeurs de texte utilisent les expressions régulières (ER) pour la recherche et le remplacement
	\item La plupart des langages de programmation introduisent des mécanismes pour utiliser les ER
	\item Extraction de données : par exemple, \expword{extraire les émails et les numéros de téléphones à partir des blogs et des réseaux sociaux}
	\item \url{https://github.com/kariminf/aruudy}
	\begin{itemize}
		\item \textit{One of my crazy projects}
		\item Détection du mètre d'une verse de poésie arabe
		\item Avantage : les règles sont lisibles
		\item Inconvénient : plusieurs passes pour traiter une verse 
	\end{itemize}
\end{itemize}

\end{frame}

\subsection{Distance d'édition}

\begin{frame}
\frametitle{Traitements basiques du texte : Caractères}
\framesubtitle{Distance d'édition}

\begin{itemize}
	\item 
\end{itemize}

\end{frame}

%===================================================================================
\section{Segmentation du texte}
%===================================================================================

\begin{frame}
\frametitle{Traitements basiques du texte}
\framesubtitle{Segmentation du texte}


\end{frame}

\subsection{Délimitation de la phrase}

\begin{frame}
\frametitle{Traitements basiques du texte : Segmentation du texte}
\framesubtitle{Délimitation de la phrase} 

\begin{itemize}
\item 
\end{itemize}

\end{frame}

\subsection{Séparation des mots}

\begin{frame}
\frametitle{Traitements basiques du texte : Segmentation du texte}
\framesubtitle{Séparation des mots}

\begin{itemize}
\item 
\end{itemize}

\end{frame}

%===================================================================================
\section{Normalisation et filtrage des mots}
%===================================================================================

\begin{frame}
\frametitle{Traitements basiques du texte}
\framesubtitle{Normalisation et filtrage des mots}


\end{frame}

\subsection{Normalisation des mots}

\begin{frame}
\frametitle{Traitements basiques du texte : Normalisation et filtrage}
\framesubtitle{Normalisation des mots}

\begin{itemize}
\item 
\end{itemize}

\end{frame}

\subsection{Filtrage des mots}

\begin{frame}
\frametitle{Traitements basiques du texte : Normalisation et filtrage}
\framesubtitle{Filtrage des mots}

\begin{itemize}
\item 
\end{itemize}

\end{frame}

%===================================================================================
\section{Lexique}
%===================================================================================

\begin{frame}
\frametitle{Traitements basiques du texte}
\framesubtitle{Lexique}


\end{frame}

\subsection{Formation des mots}

\begin{frame}
\frametitle{Traitements basiques du texte : Lexique}
\framesubtitle{Formation des mots}

\begin{itemize}
\item Flexion
\item Dérivation
\end{itemize}

\end{frame}

\subsection{Réduction des formes des mots}

\begin{frame}
\frametitle{Traitements basiques du texte : Lexique}
\framesubtitle{Réduction des formes : Racinisation}

\begin{itemize}
\item 
\end{itemize}

\end{frame}

\begin{frame}
\frametitle{Traitements basiques du texte : Lexique}
\framesubtitle{Réduction des formes : Lemmatisation}

\begin{itemize}
	\item 
\end{itemize}

\end{frame}

\insertbibliography{TALN02}{*}

\end{document}

